% Part: incompleteness
% Chapter: arithmetization-syntax
% Section: coding-terms

\documentclass[../../../include/open-logic-section]{subfiles}

\begin{document}

\olfileid{inc}{art}{trm}
\olsection{পদের সংকেতায়ন}

\begin{explain}
পদ আসলে বিশেষ ধরনের প্রতীক-অনুক্রম: পদের গঠনবিধি অনুসারে ধ্রুবক ও চলরাশি
থেকে এটি আরোহমূলকভাবে গঠিত। প্রতীকের জন্য একটি সংকেতায়ন এবং সংখ্যা-অনুক্রম
সংকেতায়নের একটি উপায় ব্যবহার করে প্রতীক-অনুক্রমকে সংখ্যা দিয়ে সংকেতায়িত
করা যায়; তাই পদকে গ্যোডেল সংখ্যা দেওয়া কঠিন নয়। আসল সমস্যা হলো, কোনো
সংখ্যা সুগঠিত পদের গ্যোডেল সংখ্যা হওয়ার ধর্মটি গণনসাধ্য—বস্তুত আদিম
পুনরাবৃত্ত—তা দেখানো।
\end{explain}

!!^{variable}s ও !!{constant}s সবচেয়ে সরল পদ, এবং $x$ এমন পদের গ্যোডেল
সংখ্যা কি না পরীক্ষা করা সহজ। কোনো~$i$-এর জন্য $x$ যদি $\Gn{\Obj v_i}$ হয়,
তবে $\fn{Var}(x)$ সত্য। অন্যভাবে বললে, $x$ দৈর্ঘ্য~$1$-এর একটি অনুক্রম এবং
তার একমাত্র উপাদান $(x)_0$ কোনো !!{variable}~$\Obj v_i$-এর কোড; অর্থাৎ
কোনো~$i$-এর জন্য $x$ হলো $\tuple{\tuple{1, i}}$। একইভাবে, কোনো~$i$-এর জন্য
$x$ যদি $\Gn{\Obj c_i}$ হয়, তবে $\fn{Const}(x)$ সত্য। এমন $i$ থাকলে তা
অবশ্যই $< x$; তাই এই দুটি সম্বন্ধই আদিম পুনরাবৃত্ত:
\begin{align*}
  \fn{Var}(x) & \defiff \bexists{i<x}{x = \tuple{\tuple{1, i}}}\\
  \fn{Const}(x) & \defiff \bexists{i<x}{x = \tuple{\tuple{2, i}}}
\end{align*}

\begin{prop}
\ollabel{prop:term-primrec}
$x$ যথাক্রমে কোনো পদ অথবা বদ্ধ পদের গ্যোডেল সংখ্যা হলে এবং কেবল হলেই সত্য
সম্বন্ধ $\fn{Term}(x)$ ও $\fn{ClTerm}(x)$ আদিম পুনরাবৃত্ত।
\end{prop}

\begin{proof}
প্রতীক-অনুক্রম~$s$ তখনই একটি পদ, যখন $s_0$, \dots, $s_{k-1} = s$ পদের এমন
একটি অনুক্রম আছে, যা পদের গঠনবিধি অনুসারে !!{constant}s ও !!{variable}s থেকে
$s$ কীভাবে গঠিত হয়েছে তা নথিভুক্ত করে। এমন সম্ভাব্য গঠন-অনুক্রম গঠনবিধি
মানছে—এ কথা প্রকাশ করতে প্রত্যেক $i < k$-এর জন্য নিচের একটির সত্য হওয়া চাই:
\begin{enumerate}
\item $s_i$ একটি !!a{variable}~$\Obj v_j$, অথবা
\item $s_i$ একটি !!a{constant}~$\Obj c_j$, অথবা
\item $s_i$ হলো $i$-তম স্থানের আগে উপস্থিত $n$টি পদ $t_1$, \dots, $t_n$
  থেকে একটি $n$-স্থানীয় !!{function}~$\Obj f^n_j$ ব্যবহার করে গঠিত।
\end{enumerate}
গ্যোডেল সংখ্যার উপর অনুরূপ সম্বন্ধটি আদিম পুনরাবৃত্ত দেখাতে এই শর্তকে আদিম
পুনরাবৃত্ত অপেক্ষক, সম্বন্ধ ও আবদ্ধ পরিমাণায়ন দিয়ে প্রকাশ করতে হবে।

ধরি, $y$ হলো $s_0$, \dots, $s_{k-1}$ অনুক্রমটির কোড; অর্থাৎ
$y = \tuple{\Gn{s_0}, \dots, \Gn{s_{k-1}}}$। এটি গ্যোডেল সংখ্যা~$x$-বিশিষ্ট
পদের একটি গঠন-অনুক্রম সংকেতায়িত করে যদি এবং কেবল যদি প্রত্যেক $i < k$-এর
জন্য:
\begin{enumerate}
\item $\fn{Var}((y)_i)$, অথবা
\item $\fn{Const}((y)_i)$, অথবা
\item এমন $n$ ও সংখ্যা~$z = \tuple{z_1, \dots, z_n}$ আছে যাতে প্রত্যেক
  $z_l$ কোনো $i' < i$-এর জন্য $(y)_{i'}$-এর সমান এবং
\[
(y)_i = \Gn{\Obj f^n_j(} \concat \fn{flatten}(z) \concat \Gn{)},
\]
\end{enumerate}
এবং আরও $(y)_{k-1} = x$। (অপেক্ষক $\fn{flatten}(z)$ অনুক্রম
$\tuple{\Gn{t_1}, \dots, \Gn{t_n}}$-কে $\Gn{t_1, \dots, t_n}$-এ পাঠায়
এবং এটি আদিম পুনরাবৃত্ত।)

(3)-এ সূচক $j$, $n$, পদ $t_l$-গুলির গ্যোডেল সংখ্যা $z_l$, এবং
$\tuple{z_1, \dots, z_n}$ অনুক্রমের কোড~$z$—সবই~$y$-এর চেয়ে ছোট। উপরের
$k$-এর জায়গায় $\len{y}$ বসানো যায়। তাই ``$y$ হলো গ্যোডেল সংখ্যা~$x$-বিশিষ্ট
পদের গঠন-অনুক্রমের কোড'' কথাটি এমনভাবে প্রকাশ করা যায়, যা সম্বন্ধটিকে আদিম
পুনরাবৃত্ত বলে দেখায়।

এখন শুধু নিশ্চিত হতে হবে যে~$y$-এর একটি আদিম পুনরাবৃত্ত সীমা আছে। কিন্তু
$x$ যদি কোনো পদের গ্যোডেল সংখ্যা হয়, তবে তার সর্বাধিক $\len{x}$টি পদবিশিষ্ট
একটি গঠন-অনুক্রম থাকবেই (কারণ~$s$-এর গঠন-অনুক্রমের প্রতিটি পদকে~$s$-এর কোনো
একটি স্থান থেকে শুরু হতে হবে, এবং দুটি উপপদ একই স্থান থেকে শুরু হতে পারে না)।
$s$-এর প্রতিটি উপপদের গ্যোডেল সংখ্যা অবশ্যই $\le x$। অতএব
$k=\len{x}$ হলে সবসময় $\le p_{k-1}^{k(x+1)}$ কোডবিশিষ্ট একটি গঠন-অনুক্রম
থাকে।

$\fn{ClTerm}$-এর জন্য !!{variable}s-এর শর্তটি শুধু বাদ দাও।
\end{proof}

\begin{prob}
দেখাও, $\tuple{\Gn{t_1}, \dots, \Gn{t_n}}$ অনুক্রমকে
$\Gn{t_1, \dots, t_n}$-এ পাঠানো অপেক্ষক $\fn{flatten}(z)$ আদিম পুনরাবৃত্ত।
\end{prob}

\begin{prop}\ollabel{prop:num-primrec}
  অপেক্ষক $\fn{num}(n) = \Gn{\num{n}}$ আদিম পুনরাবৃত্ত।
\end{prop}

\begin{proof}
  আদিম পুনরাবৃত্তি দিয়ে $\fn{num}(n)$ সংজ্ঞায়িত করি:
  \begin{align*}
    \fn{num}(0) & = \Gn{\Obj 0}\\
    \fn{num}(n+1) & = \Gn{\prime(} \concat \fn{num}(n) \concat \Gn{)}.
  \end{align*}
\end{proof}

\end{document}
